\section{Real Benchmark Cases Illustrating Why \plugmem Outperforms Task-Agnostic and Task-Specific Baselines}
\label{app:real_benchmark_cases}
\subsection{LongMemEval}

\paragraph{Case Study: Hierrachical Retrieval for Episodic Memory.}
Table~\ref{tab:case_longmemeval_2} show an examples of \plugmem retrieving episodic memory hierrachically on LongMemEval. 

\begin{table}[H]
\centering
\small
\caption{Example 1 of retrieval trace showing hierarchical retrieval for episodic memory.}
\label{tab:case_longmemeval_2}

\setlength{\tabcolsep}{8pt}
\renewcommand{\arraystretch}{1.1}
\begin{tabularx}{\linewidth}{p{0.97\linewidth}}

\toprule
\textbf{LongMemEval Case Study: Hierarchical Retrieval}\\
\midrule
\textbf{Query.} \emph{"How many weddings have I attended in this year?"}\\
\midrule
\textbf{Semantic Memory Retrieved}\\
\midrule
\(\triangleright\) \textbf{Semantic Memory 669} : \emph{User's sister's wedding was amazing, and User had a great time planning it with her, including choosing the dress and deciding on the menu.}\\
\(\triangleright\) \textbf{Semantic Memory 141} : \emph{User is planning a wedding and dreams of having a small, outdoor ceremony at a beach or in a park. User recently attended their college roommate's wedding in the city which had a rooftop garden ceremony overlooking the skyline.}\\
\(\triangleright\) \textbf{Semantic Memory 146} : \emph{Social media platforms like Instagram and Pinterest can be used to follow wedding planners, venues, and photographers. Reading reviews on Yelp or Google provides insights into the venue's reputation.}\\
...\\
\midrule
\textbf{Episodic Memory Retrieved}\\
\midrule
\(\triangleright\) \textbf{Episodic Memory Session 24} : \emph{Date: 2023/10/15 (Sun) 04:44 User: I'm planning my own wedding and I was wondering if you could give me some tips on how to choose the perfect venue. By the way, I just got back from a friend's wedding last weekend, and it was amazing - the bride, Jen, looked stunning in her bohemian-inspired dress, and her husband, Tom, was clearly smitten with her. It was at a rustic barn in the countryside, and it was so cozy and relaxed...}\\
\(\triangleright\) \textbf{Episodic Memory Session 6} : \emph{Date: 2023/10/15 (Sun) 10:57 User: I'm planning my own wedding and I need some help with finding a venue. I've always dreamed of having a small, outdoor ceremony, maybe at a beach or in a park. Do you have any recommendations or resources for finding the perfect spot? By the way, I just got back from my college roommate's wedding in the city, and it was beautiful - they had a rooftop garden ceremony overlooking the skyline, and it was so romantic. My friend Emily finally got to tie the knot with her partner Sarah, and it was amazing to see them so happy after waiting for marriage equality... }\\
...\\
\midrule
\textbf{Reasoning}\\
\midrule
\(\triangleright\) \emph{To determine how many weddings you have attended this year, let's review the information from the relevant memories provided:}\\
\emph{1. \textbf{Cousin's Wedding at a Vineyard (August)}: You mentioned attending your cousin Rachel's wedding at a vineyard in August.}\\
\emph{2. \textbf{College Roommate's Wedding in the City}: You mentioned attending your college roommate's wedding, which had a rooftop garden ceremony.}\\
\emph{3. \textbf{Friend Jen's Wedding at a Rustic Barn}: You mentioned attending your friend Jen's wedding at a rustic barn in the countryside last weekend...}\\
\midrule
\textbf{Answer.} \textbf{"Therefore, you have attended 3 weddings this year."}.\\
\bottomrule
\end{tabularx}
\end{table}
\paragraph{Takeaways}
The example highlights a key reason why \plugmem outperforms task-agnostic baselines on LongMemEval. Unlike Zep \cite{rasmussen2025zeptemporalknowledgegraph}, which contains episode nodes in its graph but only performs retrieval at the semantic level (semantic entities and facts), and LiCoMemory \cite{huang2026licomemorylightweightcognitiveagentic}, which performs graph retrieval for original dialogue chunks solely based on entities, \plugmem retrieves episodic memories by leveraging semantic memory extracted from those episodic memories. Specifically, the retrieval module first identifies semantic memories relevant to the query. In the example, semantic memory 669, semantic memroy 141, and semantic memory 146 are all related to wedding, which is the topic of the question. It then locates the source episodic memory corresponding to each retrieved semantic memory and selects those containing a sufficient number of semantic memories retrieved. In the example, semantic memory 669 and some other retrieved memories are extracted from episodic memory session 24, so episodic memory session 24 is selected. So does episodic memory session 6. We can notice that, although question-related, semantic memory 141 and 146 do give sufficient information about which wedding they are describing, making them ambiguous when counting the number. But by tracing back to its source episodic memory, \plugmem successfully distinguishes them, showing the importance of episodic memory which is neglected in Zep. Also, although LiCoMemory retrieves back origin dialogue chunks, the retrieval process is based on semantic entities. Due to the fact that episodic memory related a single entity is definitely more than episodic memory related to a specific piece of semantic memory, it is less efficiency and may bring more noise. 

\subsection{HotpotQA}
\paragraph{Case Study: Bridge-Entity Multi-hop Retrieval.}
Table~\ref{tab:case_hotpotqa_1} and Table ~\ref{tab:case_hotpotqa_2} show two examples of multi-hop retrieval on HotpotQA alternating abstract node like concept and specific node like semantic node. 

\begin{table}[ht]
\centering
\small
\caption{Example 1 of retrieval trace showing bridge-entity discovery and evidence compression.}
\label{tab:case_hotpotqa_1}

\setlength{\tabcolsep}{8pt}
\renewcommand{\arraystretch}{1.1}
\begin{tabular}{p{0.97\linewidth}}
\toprule
\textbf{HotpotQA Case Study 1: Bridge-Entity Multi-hop Retrieval}\\
\midrule
\textbf{Query.} \emph{``\;`One Less Set of Footsteps' is a song written and performed by a singer born in which year?\;''}\\
\midrule
\textbf{Hop 0 (init).} Tags for query: {[\textcolor{yes}{One Less Set of Footsteps}, \textcolor{no}{songwriter}, \textcolor{no}{performer}, \textcolor{partial}{birthdate}]}\\
\(\triangleright\) \textbf{node 5891} : \emph{``One Less Set of Footsteps'' is a song written and performed by Jim Croce. It was released in 1973 as the first single from his album ``Life and Times''.}   (\textcolor{yes}{bridge entity identified: \textbf{Jim Croce}})\\
\(\triangleright\) node 5892 : "One Less Set of Footsteps" reached a peak of \#37 on the "Billboard" Hot 100, spending ten weeks on the chart.\\
...... \\
\(\triangleright\) node 15424: Mariah Carey, born on March 27, 1969 or 1970, is an American singer, songwriter, record producer, and actress.\\
\midrule
\textbf{Query\_new} = \texttt{Integrated} ( Query, top-k nodes)\\
\midrule
\textbf{Hop 1 (refine \& expand).} Tags for query:  {[\textcolor{yes}{Jim Croce}, \textcolor{yes}{birth year}]}\\
\(\triangleright\) \textbf{node 5897} : \emph{James Joseph ``Jim'' Croce was born on January 10, 1943, and died on September 20, 1973.} (\textcolor{yes}{golden fact retrieved: \textbf{Jan 10, 1943}}) \\
node 5895 : James Joseph "Jim" Croce was an American folk and rock singer active between 1966 and 1973. He released five studio albums and singles during this period. \\
...... \\
node 5883 : "It Doesn't Have to Be That Way" was originally released early in 1973 as the B-side of the "One Less Set of Footsteps" single and was reissued in December of the same year as the third and final single from the album "Life and Times".\\
\midrule
\textbf{Reasoning (compressed).}\\
\(\triangleright\) \emph{``One Less Set of Footsteps'' is written and performed by Jim Croce.}\\
\(\triangleright\) \emph{Jim Croce was born on January 10, 1943.}  \\
\midrule
\textbf{Answer.} \textbf{1943}.\\
\bottomrule

\end{tabular}

\end{table}

\begin{table}[H]
\centering
\small
\caption{Example 2 of retrieval trace showing bridge-entity discovery and evidence compression.}
\label{tab:case_hotpotqa_2}

\setlength{\tabcolsep}{8pt}
\renewcommand{\arraystretch}{1.1}
\begin{tabular}{p{0.97\linewidth}}

\toprule
\textbf{HotpotQA Case Study 2: Bridge-Entity Multi-hop Retrieval}\\
\midrule
\textbf{Query.} \emph{``Bill Nelson flew as a Payload Specialist on a Space Shuttle launched for the first time in what year?''}\\
\midrule
\textbf{Hop 0 (init).} Tags for query:
{[\textcolor{yes}{Bill Nelson}, \textcolor{partial}{Payload Specialist}, \textcolor{partial}{Space Shuttle}, \textcolor{no}{first launch}]}\\
\(\triangleright\) \textbf{node 1375} (v=0.501):
\emph{In January 1986, Clarence William Nelson II became the first sitting member of the United States House of Representatives to fly in space, serving as a Payload Specialist on the Space Shuttle Columbia.}
(\textcolor{yes}{bridge entity identified: \textbf{Space Shuttle Columbia}})\\
\(\triangleright\) node 1371:
In 1983, Byron Kurt Lichtenberg and Ulf Merbold became the first Payload Specialists to fly on the shuttle.\\
......\\
\(\triangleright\) node 1383:
Dirk Dries David Damiaan, Viscount Frimout, flew aboard NASA Space Shuttle mission STS-45 as a payload specialist, making him the first Belgian in space.\\
\midrule
\textbf{Query\_new} = \texttt{Integrated} (Query, top-k nodes)\\
\midrule
\textbf{Hop 1 (refine \& expand).} Tags for query:
{[\textcolor{yes}{Space Shuttle Columbia}, \textcolor{yes}{first launch}]}\\
\(\triangleright\) \textbf{node 1365}:
\emph{Space Shuttle ``Columbia'' (Orbiter Vehicle Designation: OV-102) launched for the first time on mission STS-1 on April 12, 1981, marking the first flight of the Space Shuttle program.}
(\textcolor{yes}{golden fact retrieved: \textbf{1981}})\\
\(\triangleright\) node 1380:
STS-61-C's seven-person crew included \ldots\ and Representative Bill Nelson (D-FL), the second sitting politician to fly in space.\\
......\\
\(\triangleright\) node 1384:
A payload specialist (PS) is an individual selected and trained \ldots\ for a NASA Space Shuttle mission.\\
\midrule
\textbf{Reasoning (compressed).}\\
\(\triangleright\) \emph{Bill Nelson (Clarence William Nelson II) served as a Payload Specialist on the Space Shuttle Columbia.}\\
\(\triangleright\) \emph{Space Shuttle Columbia first launched on April 12, 1981 (STS-1).}\\
\midrule
\textbf{Answer.} \textbf{1981}.\\
\bottomrule
\end{tabular}

\end{table}

\paragraph{Takeaways. } 
Across both case studies, \plugmem succeeds by explicitly separating \emph{what to retrieve next} from \emph{where the evidence resides}. In Hop0, the retriever typically surfaces a small set of low-level propositions that are semantically close to the question but may only partially resolve it. Crucially, rather than repeatedly expanding within a local neighborhood around these propositions (as in many graph-based retrievers that rely on $1$--$2$ hop adjacency), \plugmem elevates the intermediate result into an \emph{abstract routing signal}, a bridge concept (e.g., \emph{Jim Croce}) or a bridge entity/context (e.g., \emph{Space Shuttle Columbia}). This abstraction step turns multi-hop QA into a sequence of targeted sub-queries: first identify the missing bridge, then retrieve the specific fact needed to answer. 


Technically, this behavior is enabled by the bipartite organization of our memory graphs, where high-level nodes (concepts/intents) connect to many low-level nodes (propositions/prescriptions).
Routing through the abstract layer allows the retriever to ``jump'' beyond the narrow vicinity of an initial seed proposition and activate a much broader set of candidate evidence, which is especially important on HotpotQA where supporting facts often lie in different articles or distant parts of the corpus. In both examples, the first hop retrieves a bridge-bearing statement; the next hop uses the bridge to directly access the answer-bearing statement (birth year or first-launch year), while irrelevant but superficially related candidates (e.g., other payload specialists, other singer-songwriters) are naturally deprioritized once the bridge entity is fixed.


Finally, the reasoning module plays a complementary role by compressing the retrieved pool into a minimal sufficient evidence set (typically two propositions in these examples), which improves robustness and token efficiency. Overall, \plugmem's abstraction-to-specificity routing and budgeted candidate control expand the reachable search space and deepen evidence chaining, without being constrained to short range neighbor expansion from an arbitrary starting node. 


\subsection{WebArena}

\paragraph{Case Study: \plugmem Dynamically Construct Useful Guidance for Agent through Retrieval and Reasoning}

This WebArena case study demonstrates that the advantage of \plugmem lies not in executing a particular workflow correctly, but in how guidance is constructed, adapted, and consumed during decision-making. As shown in Table~\ref{tab:case_webarena_1}, \plugmem enables the agent to solve the task without relying on static workflows or verbose episodic recall, instead producing compact, step-adaptive guidance that directly reflects the agent’s current informational bottleneck.\\
Compared to \textbf{task-agnostic} retrieval-based methods, the core limitation exposed by this task is not retrieval coverage but retrieval granularity. Flat retrieval systems can surface past trajectories or similar task descriptions, but these memories are typically either too specific (tied to particular pages or layouts) or too verbose to be directly actionable. As a result, the agent must internally interpret and reconcile large amounts of loosely relevant information, leading to context explosion and brittle decision making. In contrast, \plugmem retrieves procedural knowledge units—such as category navigation, constraint-based filtering, and ranking under partial observability—and further compresses them through its reasoning module into concise, decision-aligned instructions. The benefit is not simply fewer tokens, but a sharper alignment between retrieved memory and the agent’s immediate control decision.\\
On the other end of the spectrum, \textbf{task-specific} web agents such as AWM and SteP hard-code or accumulate a single canonical workflow per domain. While effective when tasks closely match the assumed structure, such workflows lack the ability to adapt their internal logic as the task unfolds. This case highlights that web navigation tasks are not monolithic: the agent alternates between qualitatively different reasoning regimes (e.g., locating the correct semantic scope, enforcing numeric constraints, and resolving ties under incomplete information). A static workflow cannot anticipate these shifts. Although AWM updates its workflow after task completion, it does not revise guidance within a task or condition it on intermediate observations. \plugmem, by contrast, recomputes guidance at each step, allowing the agent’s strategy to evolve as the environment reveals new structure.\\
More broadly, this example illustrates that \plugmem’s advantage comes from treating knowledge as the unit of memory access, rather than raw trajectories or fixed workflows. Task-agnostic methods fail because they retrieve too much unstructured experience; task-specific methods fail because they retrieve too little adaptability. \plugmem occupies a middle ground: it abstracts experience into reusable procedural knowledge while retaining the ability to contextualize and re-specialize that knowledge through reasoning. This combination allows \plugmem to generalize across WebArena tasks while still producing highly targeted, step-specific guidance, which neither task-agnostic nor task-specific designs can reliably achieve.


\begin{table}[H]
\centering
\small
\caption{Example 1 of \plugmem assisting WebArena agent.}
\label{tab:case_webarena_1}

\setlength{\tabcolsep}{8pt}
\renewcommand{\arraystretch}{1.1}
\begin{tabularx}{\linewidth}{p{0.97\linewidth}}

\toprule
\textbf{WebArena Case Study: Adaptive Retrieval and Refined Guidance}\\
\midrule
\textbf{Query.} \emph{``\;Buy the highest rated product from the meat substitute category within a budget between 100 and 200.\;''}\\
\midrule
\textbf{\plugmem Reasoning.}\\
\(\triangleright\) \textbf{Reasoning:}\\The user's goal is to buy the highest-rated product from the "meat substitute" category within a budget of 100 to 200. The current observation does not provide direct access to the "meat substitute" category or relevant products. Additionally, the visible products on the page do not appear to belong to the "meat substitute" category. To efficiently locate the desired product, the agent should use the dropdown menu to navigate to the "meat substitute" category and then apply a price filter to narrow down the options to the specified budget range. If the platform supports advanced search, it can also be used to filter by category and price range.\\
\(\triangleright\) \textbf{Final Information}:\\1. Use the dropdown menu to navigate to the "Grocery \& Gourmet Food" category, then locate the "meat substitute" subcategory.\\ 2. Once in the "meat substitute" category, apply a price filter to display products within the 100 to 200 range.\\3. Sort the results by rating to identify the highest-rated product within the budget.\\4. If the platform supports advanced search, use it to filter by category ("meat substitute") and price range (100 to 200) to streamline the process.\\5. Once the highest-rated product is identified, proceed to add it to the cart.\\
\midrule
\textbf{AgentOccam Action 1:}
\textbf{hover [1490]}. (agent operates the dropdown menu)\\
\midrule
...\\
\midrule
\textbf{\plugmem Reasoning.}\\
\(\triangleright\) \textbf{Reasoning}\\The user's goal is to buy the highest-rated product from the "Meat Substitutes" category within a budget of 100 to 200. The current page lists various products in this category, but the budget constraint narrows the options to only those priced between 100 and 200. To efficiently identify the highest-rated product within this range, the agent should apply the price filter for 100-199.99 and then sort the results by rating. This ensures that only relevant products are considered, and the highest-rated one is selected.\\
\(\triangleright\) \textbf{Final Information}\\1. Apply the price filter for the range 100.00 - 199.99 by clicking the corresponding link under the "Price" section.\\2. Once the filtered results are displayed, sort the products by "Rating" to identify the highest-rated product.\\3. Select the highest-rated product from the filtered list and click "Add to Cart."\\4. Confirm that the product meets the budget and rating criteria before proceeding.\\
\midrule
\textbf{AgentOccam Action 3:}
\textbf{click [6093]}. (agent selects the price range)\\
\midrule
\textbf{\plugmem Reasoning.}\\
\(\triangleright\) \textbf{Reasoning}\\The user's goal is to buy the highest-rated product from the "Meat Substitutes" category within a budget of 100 to 200. The current page already displays products filtered by the specified price range. The next step is to identify the product with the highest rating among the listed options and add it to the cart. The observation provides details about the products, including their names, prices, and ratings, which are essential for making the selection.\\
\(\triangleright\) \textbf{Final Information}\\
1. Review the listed products on the current page and identify the one with the highest rating. Based on the observation:\\
- "Beyond Meat Beef Beefy Crumble, 5 Pound -- 2 per case" has a rating of 50 and costs 108.27.\\
- Other products do not have visible ratings in the observation.\\
2. Since "Beyond Meat Beef Beefy Crumble" has the highest rating and falls within the budget, proceed to click the "Add to Cart" button for this product.\\
3. Confirm that the product has been successfully added to the cart before proceeding further.\\
\midrule
\textbf{AgentOccam Action 4:}
\textbf{click [10183]}. (agent inspects the item information)\\
\midrule
...\\
\bottomrule
\end{tabularx}
\end{table}

% \begin{table}[h!]
% \centering
% \renewcommand{\arraystretch}{1.2}
% \begin{tabularx}{\textwidth}{@{}c l X X@{}}
% \toprule
% \textbf{Step} & \textbf{State} & \textbf{Retrieved Memory (Procedural Summary)} & \textbf{Action \& Outcome} \\
% \midrule
% 0 &
% Home page (product showcases) &
% The target item is not visible; use the global search bar, enter the exact product name, and submit. &
% Typed full product name into the search input (UI id: 1913) and submitted. \\
% \addlinespace
% 1 &
% Search results page &
% The exact product appears in results; locate its row and click \emph{Add to Wish List} for that item. &
% Clicked \emph{Add to Wish List} on the exact-match row (UI id: 4466). Redirected to wish list. \\
% \addlinespace
% 2 &
% Wish list page &
% Verify that the target item is present; task complete. &
% Page displayed a confirmation alert that the product was added; the wish list contained the target item. Agent stopped successfully. \\
% \bottomrule
% \end{tabularx}
% \caption{Trajectory highlighting where retrieved memory directly informed the next action.}
% \label{tab:webarena_case_study}
% \end{table}